\documentclass[../../main.tex]{subfiles}
\begin{document}

{\bf[解]}

$a$を実数として，題意の直線は
\begin{align}
  2x-5y+a=0 \label{eq:1}
\end{align}
と表せる．
これと中心$(m,n)$，半径$r$の円が共有点をもつので，
$(m,n)$と直線の距離$\ell$が$r$以下だから
\begin{align}
  \ell &= \dfrac{|2m-5n+a|}{\sqrt{4+25}}\le r \\
  \therefore
  \ell &= \dfrac{|2m-5n+a|}{29} \le r \label{eq:2}
\end{align}
である．$a$を動かした時に常にこれを満たす$(m,n)$があるような$r$の最小値を求めれば良い．

ここで $2\perp 5$ だから $2m-5n$ は全ての整数値をとる．
従って $a$ の小数部分 $\alpha$ とすれば
\begin{align}
  \begin{cases}
    0\le \alpha \le \dfrac{1}{2} &\min |2m-5n+a|=\alpha \\
    \dfrac{1}{2} \le \alpha<1     &\min |2m-5n+a|=1-\alpha
  \end{cases}
\end{align}
である．従って\cref{eq:2}に代入して，いずれの場合も
\begin{align}
  \min \ell = \dfrac{1/2}{\sqrt{29}}
\end{align}
である．従って求める$r$の最小値は
\begin{align}
  \min r = \dfrac{\sqrt{29}}{58}
\end{align}
である．

\newpage
\end{document}
